% Section 2. Owner: Tejas. Sources: docs/proposal.md sections 1-2, site/index.html.
\section{Project Goals and Motivation}
\label{sec:goals}

Many computations in science and engineering are embarrassingly parallel: a large job can be cut into many pieces that never need to talk to each other. SpareCycles runs such jobs on donated machines. A client submits a job, a coordinator splits it into work units, which the design sizes at about 60 seconds each, volunteer computers pull the units, run them and return the results, and the coordinator merges the accepted results into the final output. What makes this hard is that every participant except the client may misbehave. A volunteer can disappear at any moment because its owner closes a laptop or a spot instance is reclaimed. It can return a wrong result, through faulty hardware or on purpose. And the coordinator itself can crash, taking with it the only record of which units are finished.

Most personal computers are idle for a large part of the day, and projects such as SETI@home and Folding@home have shown that this idle capacity can carry real research. Anderson estimates that buying the throughput of a typical large BOINC project on Amazon EC2 would cost about 420 times as much as running the project on volunteer machines~\cite{anderson2020boinc}. The same setting is also a demanding distributed system. The machines are anonymous and untrusted, their hardware and speed differ widely, they join and leave continuously, failures are the normal case rather than the exception, and some participants may be malicious. A platform for students and researchers must handle all of these problems at once, with a small operating budget.

We therefore set the following goals. First, a job must produce the correct result although volunteers leave or crash while they hold work. Second, it must produce the correct result although some volunteers return wrong results, by accident or by design. Third, the platform must have no single point of failure: when the coordinator that assigns work fails, another one must take over and continue the job without losing accepted work. Fourth, the system must run on AWS EC2, with three coordinators in different availability zones and volunteers on spot instances, on-demand instances and our own laptops, and every mechanism must be observable through metrics and logs. We implement all distributed algorithms ourselves.

Not every job fits this model. A suitable job splits into units that are independent, so that they can run in any order on any machine; deterministic, so that two honest volunteers return identical bytes that can be compared by hash; and compute-heavy relative to their input and output, so that shipping a unit costs far less than running it. Fractal image tiles and Monte Carlo sampling, our two workloads, have all three properties. Tightly coupled computations whose pieces exchange data, jobs that need low latency, jobs that move large data sets per unit, and computations whose results legitimately differ from one processor to another do not fit unless they are restructured.
